\section{One logical sequence, several physical forms}
\label{sec:sequence-abstractions}

The algorithms described in the preceding chapters are easiest to express over
an ordered array of nucleotide ranks. PostgreSQL does not always present a
sequence in that form. A value may be an in-memory packed object, a compressed
or out-of-line PostgreSQL value, a set of columnar fragments, or a view over a
sequence source. It may also be the result of slicing, concatenation,
reversal, or complementation. These cases have the same logical content but
different costs and lifetime constraints.

Converting every input to one byte per symbol would give algorithms a simple
interface. It would also decode complete values when a query needs only a
small interval. For a sequence of length \(N\) and a requested interval of
length \(m\), that choice changes the immediate work from \(O(m)\) to \(O(N)\).
It also creates a second representation that may be several times larger than
the packed source. This cost is particularly undesirable for repeated SQL
expressions and for alignment against long references.

BioSeqDB separates four properties of a sequence:

\begin{enumerate}
    \item The \emph{logical sequence} is a length and an ordered stream of
    alphabet ranks.
    \item The \emph{packing geometry} maps ranks to bits and physical storage
    units.
    \item The \emph{representation} owns or locates the data and defines its
    lifetime.
    \item The \emph{access requirement} states whether a consumer can stream
    symbols or requires a contiguous packed interval.
\end{enumerate}

The same logical interface can therefore serve minimizer extraction,
alignment, SQL operators, index construction, and accelerator uploads. A lazy
value may allocate metadata and retain a source handle. It is still lazy when
it postpones decoding and repacking the complete symbol payload until a
consumer requests a representation that requires them. The resident interface
is exposed by \texttt{compact\_sequence\_range}. The non-owning logical facade
is \texttt{lazy\_sequence\_view}. Figure
\ref{fig:sequence-abstraction-layers} shows this separation.

\begin{figure}[ht]
\centering
\begin{tikzpicture}[
    font=\small,
    layer/.style={draw, rounded corners, text width=11.4cm,
                  minimum height=0.78cm, align=center},
    >=Latex
]
  \node[layer,fill=blue!8] (value) at (0,3.2)
    {logical value: flat, fragmented, source-backed, sliced, concatenated, or transformed};
  \node[layer,fill=green!10] (view) at (0,2.0)
    {lazy alphabet view: length, rank stream, orientation, and optional resident-region probe};
  \node[layer,fill=orange!12] (access) at (0,0.8)
    {access decision: stream a window, borrow compatible storage, or pack an admitted interval};
  \node[layer,fill=gray!12] (range) at (0,-0.4)
    {consumer input: rank stream or packed interval};
  \node[layer,fill=purple!10] (consumer) at (0,-1.6)
    {consumer: SQL operation, minimizer, aligner, index, CPU, SIMD, or GPU path};
  \draw[->,thick] (value) -- (view);
  \draw[->,thick] (view) -- (access);
  \draw[->,thick] (access) -- (range);
  \draw[->,thick] (range) -- (consumer);
\end{tikzpicture}
\caption{Sequence abstraction layers. Ownership remains with the value or an
associated read session. Packed storage is acquired only when required by the
consumer.}
\label{fig:sequence-abstraction-layers}
\end{figure}

\section{Alphabet-generic compact ranges}
\label{sec:compact-sequence-range}

\subsection{Logical and physical traits}

The compile-time compact range represents the resident end of the design. Its
alphabet type determines rank cardinality and biological operations such as
complementation. A separate storage-unit type and bit width determine packing.
Algorithms can therefore operate on DNA4, DNA5, or another registered alphabet
without embedding a particular code table or machine-word width.

Listing \ref{lst:compact-range-interface} gives a conceptual interface. It is
an API sketch rather than an implementation listing.

\begin{lstlisting}[
  basicstyle=\ttfamily\footnotesize,
  breaklines=true,
  caption={Conceptual compact-range interface.},
  label={lst:compact-range-interface}]
compact range traits
    alphabet type
    storage unit type
    bits per symbol
    symbols per unit
    unused high bits per unit

compact range operations
    logical size
    physical unit count
    rank at logical position
    symbol iterator
    copy ranks from [first, last) into destination
\end{lstlisting}

The distinction between logical size and physical unit count is essential. If
a unit contains \(L\) symbols and the last unit contains \(P\) empty symbol
slots, then the logical length is
\[
  N = U L - P,
\]
where \(U\) is the number of units. Code that treats \(U\) as a symbol count is
correct only for unpacked storage with \(L=1\).

Packing width must also remain separate from algorithmic arithmetic. A
32-bit DNA4 storage unit holds 16 ranks, but a 27-symbol rolling key requires
more than that unit can represent. Hash state, score types, and coordinates
are selected from their mathematical ranges. They are not inferred from the
physical load width.

The range family has a common vocabulary but more than one ownership model.
Table \ref{tab:compact-range-specializations} summarizes the principal forms.

\begin{table}[ht]
\centering
\caption{Principal compact-range forms.}
\label{tab:compact-range-specializations}
\begingroup
\small
\setlength{\tabcolsep}{4pt}
\renewcommand{\arraystretch}{1.08}
\begin{tabularx}{\textwidth}{@{}
  >{\raggedright\arraybackslash}p{0.23\textwidth}
  >{\raggedright\arraybackslash}p{0.21\textwidth}
  >{\raggedright\arraybackslash}X@{}}
\textbf{Input form} & \textbf{Physical form} & \textbf{Ownership rule}\\
\hline
ordinary alphabet range & one value per symbol & borrows a stable lvalue from
the caller\\
packed owning range & bit-packed vector & owns the packed payload\\
validated PostgreSQL value & typed trailing payload & borrows the resident
datum and its established lifetime\\
\end{tabularx}
\endgroup
\end{table}

A borrowed range must not be constructed from a temporary. Likewise, a range
over a PostgreSQL value is valid only while the detoasted datum or read session
remains live. The shared interface does not erase these obligations.

\subsection{Ranks, symbols, and orientation}

Rank iteration returns compact integers suited to inner loops. Symbol
iteration converts each rank to the alphabet type and is useful when an
operation needs alphabet-aware semantics. Keeping both forms avoids repeated
symbol construction in numeric kernels while retaining a typed interface for
validation and text conversion.

Interval extraction accepts logical coordinates \([from,to)\), source
orientation, and destination order. Source orientation determines coordinate
mirroring and complementation for a reverse-complement view. Destination order
determines how the selected ranks are written. These are independent choices.
A canonical minimizer cursor, for example, can inspect reverse-complement
coordinates while retaining the order expected by its rolling state.

An allocating full conversion remains available at byte-oriented boundaries.
It is not the default for alignment. The interval form seeks to the first
relevant packed unit and decodes only the selected symbols. For resident data,
its work is \(O(to-from)\). The corresponding compact-range operation is
\texttt{to\_encoded\_into}. Lazy and fragmented representations expose the
same bounded rank-copy contract.

A logical slice has its own length but does not change the coordinate system of
the underlying payload. If a forward view covers \([a,b)\), view position
\(j\) maps to source position \(a+j\). In a reversed view it maps to
\(b-1-j\). Nested slices first normalize their bounds against the current
view, then compose this mapping with the parent. This rule avoids accumulating
rounding or padding adjustments at each SQL operation. Padding belongs to the
physical final unit, while slice bounds remain exact symbol coordinates.

Iterator categories reflect what the representation can actually provide. A
resident compact range supports direct rank access. A source-backed view may
offer the same logical operation through a bounded read rather than through a
stable reference. Generic consumers should therefore express traversal in
terms of ranks and windows. They should request a resident range only when
their kernel truly depends on typed contiguous storage.

\section{Packing and resident PostgreSQL values}
\label{sec:sequence-bit-layout}

\subsection{Packing geometry}

For an alphabet with \(\sigma\) ranks, the default bit width is
\[
  b = \lceil \log_2 \sigma \rceil.
\]
A registered packing policy may choose a wider code, provided that every rank
is representable. For a unit of \(W\) bytes, BioSeqDB stores
\[
  L = \left\lfloor \frac{8W}{b} \right\rfloor
  \quad\text{symbols and leaves}\quad
  r = 8W-Lb
\]
unused high bits. For logical position \(i\),
\[
  u(i)=\left\lfloor i/L \right\rfloor,
  \qquad s(i)=b(i\bmod L),
\]
\[
  rank(i)=\bigl(unit[u(i)]\gg s(i)\bigr)\mathbin{\&}(2^b-1).
\]
The first logical symbol occupies the least significant \(b\) bits of the
first unit. Figure \ref{fig:lsb-packed-layout} shows the layout for eight DNA4
symbols in a 16-bit unit.

\begin{figure}[ht]
\centering
\begin{tikzpicture}[x=1.05cm,y=1cm,font=\small,>=Latex]
  \foreach \x/\sym/\rank in {
    0/A/00,1/C/01,2/G/10,3/T/11,
    4/A/00,5/G/10,6/C/01,7/T/11} {
      \draw[fill=blue!8] (\x,0) rectangle ++(1,1);
      \node at (\x+0.5,0.67) {\sym};
      \node[font=\footnotesize\ttfamily] at (\x+0.5,0.28) {\rank};
  }
  \draw[->,thick] (0,-0.55) -- (7.8,-0.55)
    node[midway,below] {increasing logical position and bit significance};
  \node[anchor=east] at (-0.2,0.5) {16-bit unit};
\end{tikzpicture}
\caption{Least-significant-symbol-first packing for DNA4. Sequential
extraction masks the low rank and shifts toward higher positions.}
\label{fig:lsb-packed-layout}
\end{figure}

DNA4 uses two bits per rank and fills all standard unit widths. DNA5 needs
three bits, so each unit can contain high bits that are not symbol slots. Those
unused bits differ from the empty symbol slots in the final unit. Both values
must be recorded correctly when computing length or converting widths.

A resident rank access uses one unit load, one shift, and one mask. Repacking
an interval of \(m\) symbols takes \(O(m)\) work. Cross-width views require more
care. Matching symbol widths and divisible unit sizes are not enough to justify
a typed pointer cast. Alignment, object lifetime, permitted type access, and
byte order must also agree. BioSeqDB consequently distinguishes native typed
borrowing from byte-compatible upload. A byte-compatible descriptor may be
copied to an accelerator buffer, but it does not authorize a CPU to dereference
the bytes through another unit type.

\subsection{A flat trailing-payload representation}
\label{sec:sequence-stdc}

A C++ container is not a PostgreSQL datum. BioSeqDB therefore has a
standard-layout packed object whose fixed header is followed by aligned data.
The header records the alphabet code, storage-unit width, logical length, and
padding in the last unit. Runtime metadata is checked before it is dispatched
to an alphabet and width specialization. The inner rank loop then receives
compile-time masks, complement rules, and lane types.

The trailing payload permits a zero-copy compact range when the PostgreSQL
value is resident, uncompressed, aligned, and already validated. The adapter
does not create a vector. This statement does not mean that access to every SQL
value is copy-free. PostgreSQL may first detoast or decompress a value according
to its normal storage rules~\cite{postgresql17_toast}. It also does not mean
that a pointer remains valid after its memory context or read session ends.

Boundary validation checks header metadata, payload size, unit count, padding,
and logical length for consistency. Only after this check does runtime
metadata select a typed compact range. This places dynamic dispatch outside
the symbol loop while preserving one SQL sequence type for several alphabets
and storage widths.

\section{A lazy view over flat and fragmented storage}
\label{sec:lazy-sequence-facade}

\subsection{The reader-facing facade}

A packed pointer cannot describe a sequence that crosses fragments or draws
its symbols from a sequence source. Joining all pieces before reading the
first rank would defeat bounded access. BioSeqDB instead exposes a small
alphabet-specific view whose backing representation is selected at run time.
Its public behavior consists of:

\begin{itemize}
    \item reporting the logical length
    \item writing a requested rank interval into caller-provided storage
    \item iterating bounded chunks in either orientation
    \item probing whether a compatible packed interval is already resident
    \item reporting whether the backing state is eligible for worker sharing
\end{itemize}

The facade stores a state pointer and a table of representation operations. It
does not contain or own the complete sequence. A flat value, a fragmented
scanner, a source-backed value, and a composition tree can therefore present
the same rank-stream contract. Algorithms depend on that contract rather than
on scanner or source types.

When a requested window lies inside one fragment, a backend can read it
directly. A request that crosses boundaries is assembled in bounded scratch
storage. Reverse traversal maps the logical window to source coordinates before
changing write order and applying complement semantics. Thus a terminal slice
does not require a prefix scan through earlier fragments. Source lookup and
I/O costs still depend on the source, but the interface limits the requested
coordinate range.

The resident probe is deliberately stricter than rank streaming. It succeeds
only when the requested interval has stable storage, compatible packing, and a
valid lifetime. A source cache page is not exposed as a resident region if
the facade cannot retain its pin. This prevents a local optimization from
creating a dangling pointer.

Listing \ref{lst:lazy-window-read} summarizes a window read. The pseudocode
places coordinate normalization before backend access so that each backend
receives one valid half-open interval.

\begin{lstlisting}[
  basicstyle=\ttfamily\footnotesize,
  breaklines=true,
  caption={Bounded rank extraction through the lazy facade.},
  label={lst:lazy-window-read}]
READ_WINDOW(view, first, last, orientation, output)
    require 0 <= first <= last <= view.length
    mapped = map_to_source(first, last, orientation)
    for each intersecting piece in mapped order
        if piece is a literal
            decode only the intersecting packed ranks
        otherwise
            request only the intersecting source ranks
        apply output order and complement mapping
    require output.count = last - first
\end{lstlisting}

The output buffer is supplied by the caller, which controls temporary memory.
The backend may divide a large interval into smaller chunks, but it must
preserve rank order and report the same logical result independently of piece
boundaries. A consumer can therefore choose a chunk size suited to its cache
or vector kernel without changing sequence semantics.

\subsection{Ownership and failure handling}

The lazy facade is non-owning. A call owner or PostgreSQL read session retains
its state and releases source handles, detoasted values, arenas, and admitted
allocations. Code that invokes the facade must finish before that owner is
destroyed. Callback interfaces restrict results so that pointers, iterators,
references, and range views cannot be returned accidentally. A callback must
also avoid storing the view through a side effect.

This rule matters because PostgreSQL errors may transfer control without C++
stack unwinding. A read session registers cleanup actions and drains them on
the PostgreSQL error path before rethrowing the database error. Relying only on
ordinary destructors would leave resources live when an error crosses C++
frames~\cite{postgresql17_cpp_extensibility}.

Worker eligibility is a capability, not an ownership token. A flat resident
value may be worker-shareable. A composition is shareable only when every
source leaf permits it. Any CPU task or accelerator submission that can
outlive the current call first obtains a stable resident owner. The owner and
its memory permit remain live until task completion or the accelerator fence.

\section{Persistent slices and lazy composition}
\label{sec:sequence-piece-tree}

\subsection{Immutable piece trees}

SQL values must survive nested slicing, concatenation, and transformation. A
raw pointer into another PostgreSQL value cannot supply that lifetime. It also
cannot describe a slice that begins inside a packed unit. BioSeqDB represents
such expressions as expanded PostgreSQL objects backed by immutable balanced
piece trees.

There are three node kinds. A source leaf stores a retained source handle and
stable logical coordinates. A literal leaf stores owned bytes with validated
type information. A branch joins two subtrees and caches their combined length,
piece count, and height. Reverse and complement states are attached to edges
rather than stored as payload-bearing nodes.

Figure \ref{fig:piece-tree} illustrates the representation. Existing subtrees
can be shared by later expressions because their nodes are immutable and their
arenas retain dependencies.

\begin{figure}[ht]
\centering
\begin{tikzpicture}[
  font=\small,
  piece/.style={draw,rounded corners,align=center,text width=2.8cm,
                minimum height=0.85cm},
  >=Latex
]
  \node[piece,fill=orange!12] (root) at (0,2.7)
    {branch\\length 130};
  \node[piece,fill=blue!9] (left) at (-3.3,1.1)
    {source leaf\\\(S_1[500000,500064)\)};
  \node[piece,fill=green!10] (right) at (3.3,1.1)
    {branch\\length 66};
  \node[piece,fill=blue!9] (source) at (1.2,-0.7)
    {source leaf\\\(S_1[N-64,N)\)};
  \node[piece,fill=gray!12] (literal) at (5.0,-0.7)
    {literal leaf\\\texttt{AC}};
  \draw[->,thick] (root) -- node[left] {identity} (left);
  \draw[->,thick] (root) -- node[right] {identity} (right);
  \draw[->,thick] (right) -- node[left] {reverse} (source);
  \draw[->,thick] (right) -- node[right] {identity} (literal);
\end{tikzpicture}
\caption{Persistent composition of source slices and a literal. Reversal is
edge metadata. Neither source interval is copied into the tree.}
\label{fig:piece-tree}
\end{figure}

Concatenation joins two roots and rebalances the tree. Slicing splits at the
two boundary coordinates, retains the middle subtree, and shares unaffected
nodes. Reversal changes an edge transform and visits branch children in the
opposite logical order. At a leaf, the same transform maps coordinates,
changes output order, and complements ranks when required.

Transform composition follows the biological operation rather than the tree
shape. Two reversals cancel. Two complements cancel for alphabets whose
complement mapping is an involution. Reverse-complement can therefore be
recorded as a small edge state even when it is applied to a deeply nested
concatenation. During traversal, the accumulated state determines child order
and leaf decoding. This keeps repeated transformations independent of sequence
length and avoids inserting long chains of unary nodes.

Joins of already compatible trees and all splits follow height-bounded paths.
Concatenation that binds an unbound literal must validate and pack it. The SQL
boundary may perform additional structural and size checks.
The first expansion of a flat or out-of-line value can copy its packed bytes
once without decoding ranks. Reading a selected interval later costs \(O(m)\)
for \(m\) emitted ranks plus fragment or source overhead.

The tree's arenas make sharing compatible with PostgreSQL memory contexts.
They retain older arenas and source dependencies, and they release those
dependencies through registered cleanup. As a result, a value returned by one
SQL operator remains valid for the next operator without copying the original
symbol payload.

\subsection{Bound and unbound values}

An expanded value records whether its alphabet is bound. A bound sequence has
an alphabet and packing width. An unbound SQL literal contains canonical text
but has not adopted an alphabet. Length, slicing, reversal, comparison, and
unbound concatenation do not require complement semantics. Combining a bound
value with an unbound literal validates the literal and adopts the bound
alphabet. Two distinct bound alphabets are not silently recoded.

Element access returns a letter value containing both alphabet code and rank.
This preserves identity in arrays and prepared statements. An in-range element
of an unbound sequence is an error because its physical alphabet is unknown.
An out-of-range element yields SQL \texttt{NULL}.

Equality and hashing follow logical content. For a bound value, the key
contains a bound-state tag, alphabet code, and rank stream. Storage width,
padding, flat or expanded form, source backing, and piece boundaries do not
affect the result. An unbound value uses a different state tag and canonical
text. It is therefore not equal to a bound sequence merely because both print
the same characters.

\subsection{SQL views and transforms}
\label{sec:sequence-sql-slicing}

Sequence subscripting uses zero-based, half-open coordinates. Negative indices
are relative to logical length. Omitted bounds select the corresponding end,
and bounds are clamped to the valid range. If the normalized upper bound is
below the lower bound, the result is an empty typed sequence. Missing elements
produce \texttt{NULL}. Null bounds propagate \texttt{NULL}. Assignment,
multiple dimensions, and non-unit steps are rejected.

\begin{lstlisting}[
  basicstyle=\ttfamily\footnotesize,
  breaklines=true,
  caption={Examples of logical sequence operations.},
  label={lst:sequence-sql-examples}]
s[0]       returns the first letter
s[-1]      returns the last letter
s[:2]      selects [0, 2)
s[1:-1]    removes the first and last letters
s[4:2]     returns an empty typed sequence

concat(left, right)
reverse(sequence)
reverse_complement(bound_sequence)
\end{lstlisting}

Concatenation accepts sequence and textual operands as well as arrays of text
or letter values. Reverse is defined for bound and unbound values because it
only changes order. Reverse-complement requires a bound alphabet with a
complete complement mapping.

A cast that keeps the alphabet and storage width can remain metadata-only. A
same-alphabet width change records the target geometry and delays repacking
until a flat consumer requests it. Binding an unbound literal validates each
symbol. Relabeling a bound sequence as another alphabet fails. These rules are
the same for flat, out-of-line, expanded, and source-backed values.

\section{Access planning and materialization}
\label{sec:sequence-materialization}

\subsection{Consumer requirements}

Not every algorithm can stream. A scalar scan needs sequential ranks. A
windowed algorithm may need bounded random access. A SIMD kernel or accelerator
upload may require one contiguous packed interval. BioSeqDB represents this
choice with three access categories. Their resulting routes are summarized in
Table \ref{tab:lazy-access-routes}.

\begin{table}[ht]
\centering
\caption{Access routes selected from consumer requirements.}
\label{tab:lazy-access-routes}
\begingroup
\small
\setlength{\tabcolsep}{4pt}
\renewcommand{\arraystretch}{1.08}
\begin{tabularx}{\textwidth}{@{}
  >{\raggedright\arraybackslash}p{0.22\textwidth}
  >{\raggedright\arraybackslash}p{0.25\textwidth}
  >{\raggedright\arraybackslash}X@{}}
\textbf{Route} & \textbf{Typical case} & \textbf{Copy and lifetime rule}\\
\hline
stream symbols & forward scan over a tree, source, or fragment set & decodes
bounded chunks into caller scratch with no complete packed allocation\\
borrow resident & aligned interval in a compatible flat value & performs no
payload copy while the datum owner remains live\\
materialize admitted & unaligned interval or resident-only consumer & packs
only the selected interval and retains its allocation permit through use\\
\end{tabularx}
\endgroup
\end{table}

For sequential access, the consumer supplies bounded scratch space and reads
chunks from the facade. For contiguous access, the planner first asks whether
the representation can lend a compatible resident interval. The interval must
begin at a boundary represented by the receiving range. A start in the middle
of a packed unit cannot be passed as an ordinary compact range because the
consumer would interpret the wrong first symbol.

If borrowing is impossible, the planner computes the packed size of the
selected interval and obtains a memory permit. Materialization then packs only
\([from,to)\). The returned owner holds both the allocation and its permit.
Moving the owner into asynchronous work also moves the accounting lifetime.

The term zero-copy therefore applies to specific transitions, not to the
sequence system as a whole. A validated flat value can be adapted by reference.
A slice or transformation of an expanded tree can share payload nodes. A
fragmented interval can be streamed without a complete packed result. In
contrast, PostgreSQL may detoast a value, an unaligned interval may require
repacking, and text or binary output must create its destination form.

Durable storage is another explicit boundary. A source view is scoped to its
transaction and cannot be placed directly into a heap or index tuple. A tree
must be flattened before it crosses that boundary. If its logical size exceeds
the flat PostgreSQL representation, the operation reports a size error rather
than truncating the value.

Figure \ref{fig:lazy-lifecycle} follows the access decision from a bound SQL
expression to a consumer.

\begin{figure}[ht]
\centering
\begin{tikzpicture}[
  font=\small,
  stage/.style={draw,rounded corners,align=center,text width=3.0cm,
                minimum height=0.9cm},
  >=Latex
]
  \node[stage,fill=blue!8] (expr) at (0,3.2)
    {bound slice, concatenation, or transform};
  \node[stage,fill=orange!10] (tree) at (0,1.8)
    {expanded piece tree};
  \node[stage,fill=green!10] (view) at (0,0.4)
    {lazy alphabet view};
  \node[stage,fill=gray!10] (stream) at (-4.2,-1.4)
    {stream ranks\\bounded scratch};
  \node[stage,fill=gray!10] (borrow) at (0,-1.4)
    {borrow units\\resident and aligned};
  \node[stage,fill=gray!10] (pack) at (4.2,-1.4)
    {pack interval\\permit-owned};
  \node[stage,fill=purple!9,text width=5.2cm] (flat) at (0,-3.2)
    {flat boundary for durable storage or external representation};
  \draw[->,thick] (expr) -- (tree);
  \draw[->,thick] (tree) -- (view);
  \draw[->,thick] (view) -- node[left] {sequential} (stream);
  \draw[->,thick] (view) -- node[right] {compatible} (borrow);
  \draw[->,thick] (view) -- node[right] {fallback} (pack);
  \draw[->,thick,dashed] (stream) -- (flat);
  \draw[->,thick,dashed] (borrow) -- (flat);
  \draw[->,thick,dashed] (pack) -- (flat);
\end{tikzpicture}
\caption{A bound lazy expression can finish through any access route. Flattening is
performed only when the receiving boundary requires a durable or alternate
representation.}
\label{fig:lazy-lifecycle}
\end{figure}

\subsection{Why eager conversion is avoided}

Eager conversion has three effects beyond its immediate copy. First, it makes
small slices pay for the full source. Second, it expands compact DNA ranks into
byte or text buffers that place additional pressure on memory bandwidth.
Third, it hides allocation inside a convenience operation, which makes memory
admission and asynchronous lifetime difficult to account for.

Lazy composition changes the unit of work. Metadata operations preserve source
coordinates and transforms. Streaming consumers request bounded chunks.
Resident-only consumers name the interval they need and receive an owner for
that interval. Conversions for external byte-oriented APIs remain valid, but
they occur at a declared boundary and can be moved outside repeated calls when
their input is invariant.

This design does not assume that laziness is always faster. A consumer that
reads the complete value repeatedly may benefit from one admitted packed copy.
The access planner makes that choice visible. The important property is that a
slice or source-backed value is not forced into a complete resident vector
before the actual consumer requirement is known.

\section{Contract for sequence consumers}
\label{sec:sequence-author-contract}

The sequence model can be summarized by the following rules:

\begin{enumerate}
    \item Write alphabet-generic algorithms against compact-range traits or
    the lazy rank facade. Derive rank masks and complement behavior from the
    selected alphabet.
    \item Keep storage width separate from hash, score, and coordinate width.
    Physical packing does not define an algorithm's numeric range.
    \item Request the weakest access form that supports the computation. A
    forward scan should stream. A resident-only kernel should request and own
    the selected packed interval.
    \item Convert only the required logical interval at a byte-oriented or
    durable boundary. Do not hide complete conversion in an iterator or
    convenience wrapper.
    \item Treat every borrowed pointer as a capability tied to a named owner.
    Do not persist pointers into PostgreSQL values, source cache pages, or scratch
    buffers.
    \item Preserve logical equality across flat, fragmented, source-backed,
    and composed forms. Representation details must not change rank content.
\end{enumerate}

The earlier resident-vector model made one algorithm-friendly layout mandatory
for every input. The current model makes logical ranks and ownership mandatory,
while packed contiguity is an explicit request. This permits database-resident
slices and fragmented sources to enter the alignment pipeline without treating
all PostgreSQL and source storage as one kind of memory. The next chapter
applies these representation and lifetime rules to PostgreSQL planning and
custom scan execution.
